\documentclass[11pt]{article}
\usepackage[T1]{fontenc}
\usepackage{lmodern}
\usepackage[margin=1in]{geometry}
\usepackage{microtype,amsmath,amssymb,amsthm,mathtools}
\newcommand{\llbracket}{\mathopen{[\![}}
\newcommand{\rrbracket}{\mathclose{]\!]}}
\usepackage{booktabs,longtable,tabularx,array,enumitem}
\usepackage[dvipsnames]{xcolor}
\usepackage{hyperref}
\hypersetup{colorlinks=true,linkcolor=MidnightBlue,citecolor=MidnightBlue,urlcolor=MidnightBlue,
 pdftitle={A 238-state binary Turing machine for ZFC consistency},
 pdfauthor={Research prepared for Vladimir Reshetnikov and ProveIt}}
\setlength{\emergencystretch}{2em}
\setlength{\parskip}{0.35em}
\setlength{\parindent}{0pt}
\setlist{itemsep=0.2em,topsep=0.35em}
\newtheorem{theorem}{Theorem}[section]
\newtheorem{lemma}[theorem]{Lemma}
\newtheorem{proposition}[theorem]{Proposition}
\theoremstyle{definition}\newtheorem{definition}[theorem]{Definition}
\newcommand{\Con}{\operatorname{Con}}
\newcommand{\Halt}{\operatorname{Halt}}
\newcommand{\RHalt}{\operatorname{RegisterHalts}}
\newcommand{\loc}{\operatorname{loc}}
\newcommand{\inc}{\operatorname{inc}}
\newcommand{\strip}{\operatorname{strip}}
\newcommand{\brk}{\operatorname{brk}}
\newcommand{\cont}{\operatorname{cont}}
\newcommand{\ZF}{\mathrm{ZF}}
\newcommand{\ZFC}{\mathrm{ZFC}}
\newcommand{\PA}{\mathrm{PA}}
\newcommand{\bits}{\{0,1\}}
\newcommand{\file}[1]{\nolinkurl{#1}}
\newcommand{\digest}[1]{\par\begingroup\small\ttfamily\detokenize{#1}\par\endgroup}

\begin{document}
\begin{center}
{\Large\bfseries A 238-state binary Turing machine\\[4pt] for ZFC consistency}\par
\vspace{0.45em}
{\large Composing source regrouping with delayed erasure}\par
\vspace{0.35em}
Research prepared for Vladimir Reshetnikov and the ProveIt project\\
19 October 2026
\end{center}

\begin{abstract}
We construct an explicit binary, single-tape Turing machine with
\textbf{238 working states}, including its entire blank-tape initialization,
whose nonhalting is provably equivalent in ordinary ZFC to
$\Con(\ZFC)$.  The construction combines two distinct 239-state reductions
already present in ProveIt's incoming research: a source-boundary
reassociation with a decoder-context split, and a delayed-erasure reduction.
Their syntactically separate edits are not enough to justify composition.
We check the complete local dependencies, reprove the remaining selector
guard for the changed source, and independently verify every instruction
interface directly on the final 238-state table.

The final direct proof covers all 8,386 guarded cases of the 6,455-leaf source,
with arbitrary natural register values, arbitrary additional header gap, and
an arbitrary nonempty finite spare list.  It uses 2,255,513 symbolic proof
steps and has minimum instruction duration 13.  Independent sparse and dense
blank simulations agree on the entire tape after 409,835 transitions.
The archive contains the exact machine, source data, LaTeX, PDF, fresh
certificates, and portable standard-library Python replay.  The new endpoint
is an ordinary mathematical proof supported by executable certificates;
we do not supply a newly compiled Lean or Rocq theorem for the exact
238-state table or a serialized internal PA derivation.
\end{abstract}

\tableofcontents
\newpage

\section{Result, convention, and evidence level}
\label{sec:result}

\begin{theorem}[The exact 238-state endpoint]
Let $M_{238}$ be the transition table in
\file{machines/zfc_composed238.tm}, started in state $0$ at head position $0$
on the all-zero tape indexed by $\mathbb Z$.  State $238$ is the separate
halt outcome.  For a customary effective first-order ZFC proof predicate,
\[
 \ZFC\vdash
 \bigl(\Halt(M_{238},0^{\mathbb Z})
 \longleftrightarrow
 \exists p\,\operatorname{Proof}_{\ZFC}(p,\bot)\bigr).
 \tag{1}
\]
Consequently,
\[
 \boxed{\quad\ZFC\vdash
 \bigl(\neg\Halt(M_{238},0^{\mathbb Z})
 \longleftrightarrow\Con(\ZFC)\bigr).\quad}
 \tag{2}
\]
The theorem uses the explicitly identified inherited source-to-consistency
calibration in Section~\ref{sec:calibration}, and the new exact-table
construction and simulation proofs below.  It makes no assumption of
$\Con(\ZFC)$ in establishing the biconditionals.
\end{theorem}

The following table fixes all state-counting conventions.  In particular,
neither a supplied input word nor an uncharged initialization program is
hidden outside the machine.

\begin{center}
\begin{tabular}{ll}
\toprule
Working states & $0,1,\ldots,237$: exactly 238\\
Transition rows & 476: exactly two per working state\\
Tape and alphabet & One two-way infinite tape, alphabet $\bits$\\
Initial tape and head & All cells $0$; head at coordinate $0$\\
Start and dispatcher & $0$ and $84$\\
Separate halt outcome & $238$, with no outgoing transitions\\
Moves & Exactly one cell left or right; no stationary moves\\
In total, including halt & 239 control labels\\
\bottomrule
\end{tabular}
\end{center}

The convention is the usual working-state convention used by the supplied
240-state baseline and the cited 254-state comparison.  Thus the requested
strict threshold is met: $238<239$ working states.  Every working state is
reachable from the start in the directed transition graph.  This census
does not assert that all states occur in the actual blank run.

The canonical table has SHA-256
\digest{909f29c493e0ccc5d6a481fd7a20a9a82eb934ab2aa7df17b5ec3789255a8c08}
The inherited-label version has SHA-256
\digest{069942be27db9736875ea042320e51ac3b1eac7116cbe01cdce520fa6056f78c}
It uses start $95$, dispatcher $96$, and halt $254$.  All local row numbers
below refer to these inherited labels, unless explicitly called canonical.
The complete canonical table is also printed in Appendix~\ref{app:table}.

\subsection{What is newly established}

The inspected repository checkpoint is
\file{ce192939372ac7944fe2a0208c68a96d49fb5e3b}.  Its linked candidate remains
the 240-state source-rotation construction.  The incoming archives at that
checkpoint contain two different 239-state constructions.  Two archives
give the same delayed-erasure machine; the third gives a different machine
based on source regrouping and an entry-context split
\cite{incomingA,incomingB,incomingC,rotation240}.

The present contribution is the concrete 238-state composition and its
complete compatibility and direct-runtime verification.  The individual
239-state reductions, their source transformations, the underlying proof
enumerator and the symbolic proof rules are inherited.  We make no claim
of historical priority for those methods, of a minimum possible state
count, or of an exhaustive smallest-known-machine result.  A supplementary
formalism investigation gives a concrete obstruction to an attractive
logical-basis substitution; its precise scope is stated separately.

\begin{center}
\begin{tabularx}{\textwidth}{@{}lX@{}}
\toprule
Evidence & Status in this report\\
\midrule
Exact table, census, relabeling & Freshly reconstructed by two finite recipes.\\
Source equivalence & Complete 13-stage certificate freshly replayed.\\
Calibrated source identity & All 6,455 ordered triples and the seed freshly extracted as data from pinned Lean definitions.\\
Final binary runtime & Complete unbounded proof replay on the exact 238-state table.\\
Blank initialization & Fresh sparse/canonical and dense/inherited whole-tape agreement.\\
Local 239-to-238 proof & Six exact blocks, inductive entry facts, and a fresh predecessor guard proof.\\
Inherited logical calibration & The cited ProveIt mathematical development; its full proof-assistant ancestor library is not rebuilt here.\\
New exact-table kernel endpoint & Not supplied.  An ordinary ZFC proof and executable certificates are supplied.\\
\bottomrule
\end{tabularx}
\end{center}

\section{The two savings and their exact composition}
\label{sec:construction}

\subsection{Two different 239-state parents}

The delayed-erasure parent keeps the original 6,455-leaf source.  It removes
old control state $169$ by making states $80$ and $91$ cooperate with
existing helpers $167$ and $151$.  Its canonical table digest begins
\file{4447a2423907}; the BusyBeaver239 and Companion submissions describe
this same machine \cite{incomingA,incomingB}.

The other parent, which we denote $B_{239}$, first reparenthesizes the source
tree without changing its ordered instructions or complete control graph.
After repairing the changed decoder context, it splits the two entry roles
of old state $218$ according to the bit immediately left of the head.  One
role shares state $179$; the other occupies state $223$, formerly an internal
helper.  That construction removes $218$ and has canonical digest beginning
\file{666f5bf1370b}.  Its exact inherited-label digest is
\digest{d061770f487937f9da652919c9bc9bb79bc9bf137ca41dd7a08eb5ecda8e7719}
This table is included as \file{reference/boundary239_old.tm}.
The corresponding source has digest
\digest{7fd152d1f5371d2a561d4ad1870d241683b010ce7b7a27399b95363a4537cfb6}

The decisive observation is that $B_{239}$ still contains the precise
control fragments needed by the delayed-erasure reduction.  This is a
compatibility fact to prove, rather than a subtraction of two unrelated
state-count claims.

\subsection{The full final edit}

Starting with the exact $B_{239}$ above, perform precisely the three retained
row changes in Table~\ref{tab:edits} and delete both rows of state $169$.
All remaining retained rows are unchanged.

\begin{table}[htbp]
\centering
\begin{tabular}{rrll}
\toprule
State & Read & $B_{239}$ transition & $M_{238}$ transition\\
\midrule
147 & 0 & $(0,R,169)$ & $(0,R,80)$\\
80 & 0 & $(0,R,25)$ & $(0,L,106)$\\
91 & 1 & $(0,L,50)$ & $(1,L,167)$\\
\bottomrule
\end{tabular}
\caption{Every retained-row edit for the final 239-to-238 step.}
\label{tab:edits}
\end{table}

The deleted rows are $169/0=(0,R,57)$ and $169/1=(0,R,56)$.
The complete incoming-edge census shows that $147/0$ is the only edge into
$169$.  Thus no dangling transition remains after redirection and deletion.
There are $478-2=476$ rows; three retained rows change and 473 retained rows
remain identical.

Canonical numbering puts inherited state $95$ first, followed by every
other surviving working label in increasing order.  Inherited halt $254$
maps to canonical halt $238$.  This is an action-preserving bijection;
it sends dispatcher $96$ to $84$ and preserves every step exactly.

For incorporation against the older linked 240-state table, the archive
also supplies the complete alternate recipe
\file{reference/row_edits_from_rotation240.json}.  It removes states $169$
and $218$ and changes 43 retained rows.  The builder independently performs
this full patch and the three-row construction above, and requires
byte-identical final inherited and canonical tables.  The larger patch
includes the inherited source-routing changes; those changes must not be
mistaken for a free two-state deletion from the original table.

\subsection{Why syntactic noninterference is insufficient}

The local delayed-erasure proof uses an invariant about nearby tape cells.
Changing the source decoder can change the tapes on which those local rows
are entered, even when the rows themselves are untouched.  Therefore the
following reasoning would be invalid: the edit sets are disjoint, so both
savings survive.  We instead establish three concrete facts:

\begin{enumerate}
\item Every row needed by the local block identities is still present with
the required action.  The complete incoming-edge implications for the
structural invariant are checked on $B_{239}$ itself.
\item The only remaining semantic guard is proved anew for the changed
source and the exact $B_{239}$ initialization and runtime.
\item Every source-instruction interface is independently checked on the
final $M_{238}$ table, giving a direct correctness route.
\end{enumerate}

This combination is useful beyond the particular number 238.  A local
proof certificate has a finite dependency set of transition rows.  If those
rows remain identical, its exact local equations remain valid.  Reachability
premises are a different dependency and must be preserved by a proved
invariant or re-established on the new source.  Explicitly keeping these
two kinds of dependency separate makes compositions reviewable.

\section{The short relational proof of the final saving}
\label{sec:local}

Write a configuration as $(q,h,T)$, where $T:\mathbb Z\to\bits$.
Let $\phi(169)=80$, with $\phi(q)=q$ at every other inherited label.
At simulation checkpoints, the old and new heads and entire tapes agree;
the controls are related by $\phi$.

\subsection{The structural invariant}

The following predicate is inductive for $B_{239}$ from state $95$, even
on arbitrary initial binary tapes:
\[
\begin{array}{c|l}
q & \text{required nearby bits}\\\hline
67 & T(h-1)=0\\
80 & T(h-2)=0,\quad T(h-1)=1\\
91 & T(h-3)=0,\quad T(h-2)=1,\quad T(h-1)=1\\
169& T(h-1)=0.
\end{array}
\tag{3}
\]
There is no requirement at any other state.  The initial state is $95$,
so the predicate initially holds vacuously.  For every transition into one
of the four displayed states, the certificate checks the required shifted
bits against the bit just written and the predecessor's invariant.
All incoming transitions, including every row outside this small fragment,
are inspected.  This prevents a missed caller from invalidating the proof.

One extra premise is needed:
\[
 q=169\ \wedge\ T(h)=1\quad\Longrightarrow\quad T(h+1)=1.
 \tag{4}
\]
Unlike (3), we do not claim (4) for arbitrary binary initial tapes.  We
prove it for the blank run and the valid runtime frames of the changed
source in Section~\ref{sec:guard}.

\subsection{Six exact whole-tape equations}

Table~\ref{tab:blocks} lists all exceptional blocks.  Unspecified tape cells
are arbitrary and independent.  The local verifier rejects any read of an
unspecified exterior cell, and compares every modified cell, final head,
and endpoint control.

\begin{table}[htbp]
\centering\small
\begin{tabular}{rllrrrl}
\toprule
Old $q$ & Fixed relative bits & New $q$ & Old steps & New steps & $\Delta h$ & Endpoint\\
\midrule
80  & $T(-1)=1,T(0)=0$ &80&1&5&$+1$&25\\
80  & $T(0)=1$ &80&1&1&$+1$&91\\
169 & $T(-1)=0,T(0)=0$ &80&1&3&$+1$&57\\
169 & $T(-1)=0,T(0)=T(1)=1$ &80&1&5&$+1$&56\\
91  & $T(0)=0$ &91&1&1&$+1$&106\\
91  & $T(-3)=0,T(-2)=T(-1)=T(0)=1$ &91&3&11&$-3$&2\\
\bottomrule
\end{tabular}
\caption{The six positive-block identities.  The starting head is written
as relative coordinate zero.  Old and new final tapes agree exactly.}
\label{tab:blocks}
\end{table}

The fourth case shows the main mechanism.  The old machine erases the
scanned $1$ and enters state $56$ one cell to the right.  The new control
path is
\[
 80\longrightarrow91\longrightarrow167\longrightarrow151
 \longrightarrow173\longrightarrow56.
\]
It first retains the scanned $1$.  The guaranteed next $1$ makes the
$91/1$ action available.  Returning left, states $167$ and $151$ use the
preceding $0$ to choose the appropriate continuation.  State $173$ finally
erases the original cell and reaches exactly the old endpoint.  The
temporary retention of the bit stores a control distinction that would
otherwise require a separate working state.

For the last case, the old path $91,50,20$ erases the three consecutive
ones at relative positions $0,-1,-2$ and reaches state $2$ at position $-3$.
The new path is
\[
91,167,151,24,36,54,71,82,94,50,20.
\]
State $151$ reads $T(-2)=1$ to choose the original state-$91$ route;
it encounters $0$ in the newly shared state-$169$ route.  The additional
bit $T(-3)=0$ then licenses the $24/0$ step that begins the return excursion.
The eleven-step path produces the same three
erasures and no other tape change.  The remaining cases are identity
excursions or unchanged transitions; their full traces are in
\file{evidence/local238_certificate.json}.

\begin{lemma}[Positive-block transfer]
Assume the old execution begins at state $95$ and satisfies (4) at every
reached old state-$169$ configuration.  Then $B_{239}$ and $M_{238}$,
started on that same tape, halt together.
\end{lemma}

\begin{proof}
The initial configurations are related by $\phi$.  At an old checkpoint in
states $80,91,169$, choose the applicable row of Table~\ref{tab:blocks}.
The structural invariant (3), and (4) in its one required case, supply all
fixed bits.  The corresponding old and new positive blocks reach related
configurations with identical heads and entire tapes.  At every other
state, the old transition and the new transition agree under $\phi$;
the verifier checks all 472 such ordinary rows.  This count includes
$147/0$, which agrees after applying $\phi$.

No exceptional block halts internally.  Induction produces corresponding
checkpoints for every finite number of blocks.  Every old block takes
between one and three steps, and every new block takes between one and
eleven.  Thus an infinite old execution gives an infinite new execution
with unbounded accumulated time.  A finite old halting execution gives a
finite new halting execution.  Determinism excludes a different, earlier
new halting trajectory, yielding the converse as well.
\end{proof}

\subsection{Reproving the selector guard for the changed source}
\label{sec:guard}

The guard verifier does not reuse the old rotation240 selector theorem
merely because the local transition rows agree.  It instruments the exact
$B_{239}$ table while replaying every arbitrary-value instruction case of
the boundary source.  At each state-$169$ entry reading $1$, the symbolic
right-hand run of ones is $G+K$, with $G\ge0$.  There are 32 such symbolic
events and the minimum constant is $K=23$.  In particular its length is
at least two, which proves (4).

The bootstrap is separately checked for all 409,691 old transitions,
including six state-$169$/read-$1$ visits, and its complete 351-one-cell
endpoint is compared.  The subsequent 8,386 cases use 2,247,513 symbolic
proof steps.  Unary shortcuts only iterate a literal self-loop; state
$169$ has no self-loop.  Every one-block traversal schema is independently
derived without a visit to $169$.  Consequently no unary or list shortcut
can hide an unexamined visit relevant to (4).

The source-frame invariant and the source joining result proved below
allow this per-instruction result to be iterated indefinitely.  Therefore
(4) holds throughout the old blank run, and the local transfer lemma gives
the required blank-tape equivalence $B_{239}\leftrightarrow M_{238}$.
This relational route is supplementary: the next sections verify the final
table directly and do not require (4) as a premise.

\section{Source equivalence and the exact calibrated source}
\label{sec:source}

\subsection{Why source associativity must be checked}

The program is represented by a full binary tree whose leaves carry
increment $\mathrm{CI}(r)$, decrement-and-test $\mathrm{CD}(r)$,
break $\mathrm{CB}(k)$, continue $\mathrm{CC}(k)$, or halt
$\mathrm{CH}$.  The binary path to an instruction participates in the
control operation.  Thus an arbitrary reassociation of a source tree can
change a jump target even when the left-to-right instruction sequence is
unchanged.  The inherited boundary construction certifies the full control
behavior, including rotations across an explicit trailing break.

Number the leaf occurrences from left to right.  If $p_i$ is the path to
leaf $i$, let $\loc_T(P)$ be the leaf reached by following the infinitely
zero-padded word $P0^\omega$ in tree $T$.  The walk stops at the encountered
leaf.  Fixed-length binary increment wraps $1^d$ to $0^d$, and
$\inc(\epsilon)=\epsilon$.  Define $\strip_b(k,p)$ by deleting the suffix
through the $k$th occurrence of $b$ from the right; exhaustion yields the
empty path.  Set
\[
 \brk(k,p)=\inc(\strip_0(k,p)),\qquad
 \cont(k,p)=\strip_1(k,p).
\]
An instruction acts on the path of the leaf actually encountered, rather
than on an arbitrary unconsumed suffix of the starting program counter.
These conventions match the pinned \file{RegisterPath.lean} semantics.

The complete descriptor for a leaf records its register effect and every
successor leaf index.  Increment has one successor, decrement has both the
positive and zero successors, break and continue have their respective
jump successor, and halt has a halt descriptor.  For example,
\[
 G_T(i)=\bigl(\mathrm{CD},r,
 \loc_T(\inc(p_i)),\loc_T(\brk(1,p_i))\bigr)
\]
at a decrement leaf.  This finite array does not approximate the registers
or infer any branch to be unreachable.

\begin{lemma}[Complete descriptor equivalence]
If two finite full source trees have identical complete descriptor arrays,
then their register executions from the empty path are strongly bisimilar,
for every initial register vector.
\end{lemma}

\begin{proof}
Relate configurations whose register vectors are equal and whose decoded
leaf indices are equal.  Both empty paths decode to leaf zero.  At a matched
increment, both executions increment the same register and reach the same
next index.  At a matched decrement, equal values choose the same zero or
positive branch and give equal resulting values.  At a matched jump the
registers are unchanged and the target indices agree.  A halt descriptor
halts both sides.  Induction establishes the relation at every source
step.  Full finite source trees cannot become stuck in an internal cycle
or a missing branch.
\end{proof}

\subsection{The replayed finite certificate}

The exact inherited source transform has 13 stages and 25 tree rotations.
At every stage, the certificate checks the ordered instruction word and
all 8,385 guarded successor edges together with the unique halt descriptor.
It also requires that every unused successor suffix after the encountered
leaf is all zero.  This supplies the tape identity
\[
 P0^\omega=p_{\loc_T(P)}0^\omega,
 \tag{5}
\]
which is needed when joining binary instruction frames.  Equality of leaf
indices alone would not justify discarding an arbitrary nonzero tail.

\begin{center}
\begin{tabular}{lrr}
\toprule
Source metric & Rotation240 source $P_0$ & Boundary source $P_1$\\
\midrule
Expanded leaf occurrences & 6,455 & 6,455\\
Distinct internal subtrees & 202 & 201\\
Distinct leaf instructions & 23 & 23\\
Total distinct source DAG nodes & 225 & 224\\
Maximum leaf-path length & 51 & 48\\
Maximum consecutive zero bits & 6 & 6\\
Decrement occurrences & 1,931 & 1,931\\
Guarded instruction cases & 8,386 & 8,386\\
\bottomrule
\end{tabular}
\end{center}

The DAG-node count is a source metric, not the number of working states of
the final binary machine.  The source instruction inventory is 2,596
increments, 1,931 decrements, 128 breaks, 1,799 continues and one halt.
No instruction operand changes.  The normalized control graph has digest
\digest{ae7f12f37306b86e5f70469210fd2a8d68b6f9df094bf659abc6b4a56714f1fd}

The archive includes the original and changed leaf arrays, all rotation
paths, a compact source DAG, and the standalone source checker.  The fresh
receipt is \file{evidence/source_rotation_fresh.json}.

\subsection{Identification with the inherited logical program}
\label{sec:calibration}

The exact original source $P_0$ has digest
\digest{74b1e7f73cc3ffe7761e36ab2d57a26e0ba34d3ac624f2472799e6ac41d5a51d}
The source-interface checker parses the actual pinned Lean lookup rows
and overrides from \file{ProgramRouting245.lean},
\file{ProgramRotation244.lean} and \file{SourceRotation240.lean}.
All 203 lookup indices are checked; expanding the reachable data gives
exactly the ordered 6,455 triples of $P_0$.  This is a data-identification
check, not a fresh invocation of the Lean kernel.

The pinned path walker uses a budget of $\mathrm{nnodes}+1=204$, exceeding
the maximum source depth 51.  All reached internal references descend
strictly in node index.  Its halting predicate requires a literal halt
leaf, so neither decoder-budget exhaustion nor an internal cycle supplies
an extra halting behavior in this comparison.

The same procedure extracts the logical initial state
\[
 R_*=(3,0,0,0,3,0,0,0,0),
 \tag{6}
\]
with later logical registers zero, from \file{LogicalSeed246.lean} and
\file{L2Lib254.lean}.  The earlier calibrated theorem,
\file{CalibratedRotation240.halts_iff_zfc_refutation}, states
\begin{equation*}
 \RHalt(P_0,(\epsilon,R_*))
 \longleftrightarrow
 \operatorname{BProv}(\mathrm{ZFCax\_s},[],\bot).
 \tag{7}
\end{equation*}
The concrete inherited calibration checkpoint is
\file{13ec39ed43148b1807f51ed6e6528e1c84659562}
\cite{calibrated,source240,seed246}.

Mathematically, this calibration combines the proof enumerator, its compact
first-order calculus, and a native nonextensional set theory built from
container, global-image collection and irreflexivity axioms.  The
comparison with ZF uses the appropriate bisimulation quotient and its
internally well-founded part.  The ZF-to-ZFC refutation comparison uses the
constructible submodel and first-order soundness and completeness
\cite{wf,choice,upstream}.  These are particular ordinary mathematical
arguments that ZFC formalizes.  We do not assume a general soundness
theorem for the whole Lean type theory inside ZFC.

The endpoint $\operatorname{BProv}$ means that a finite derivation exists.
Encoding and decoding finite syntax and derivation trees identifies it
with $\exists p\,\operatorname{Proof}_{\ZFC}(p,\bot)$.  Changing between
customary effective first-order proof codings gives the same consistency
statement through the usual syntactic proof translations.  The argument
does not rest on an unqualified assertion of external equiconsistency.

Combining (7) with the source descriptor lemma gives
\[
 \RHalt(P_1,(\epsilon,R_*))
 \longleftrightarrow\neg\Con(\ZFC)
 \quad\text{provably in ZFC.}
 \tag{8}
\]
The next section identifies this exact source execution with the final
binary table directly.

\section{Direct unbounded runtime proof for the 238-state table}
\label{sec:runtime}

\subsection{Frames and their quantified parameters}

Logical registers occupy nine unary-coded physical positions, according to
the permutation
\[
 \pi=(6,2,0,4,3,7,8,5,1).
\]
Thus logical $R_i$ is stored in physical position $\pi(i)$.
Write
\[
 B(v)=0\,1^{v+1},\qquad C(v)=1\,0\,1^v.
\]
At an instruction checkpoint, the head is at coordinate $h$ in dispatcher
state $96$.  A marker $1$ occupies $h-1$.  The program-counter bits start at
$h$ and are padded with zeros through a header $1$ at $h+64+G$, where
$G\ge0$.  Immediately after the header are the nine physical register
blocks followed by a nonempty finite spare list.  The tape exterior on
both sides is zero.

For a leaf path $p$, a register vector $R\in\mathbb N^9$, a nonempty list
$A\in\mathbb N^{<\omega}$, and $G\in\mathbb N$, denote this exact whole
tape and control configuration by $F(p,R,A,G,h)$.  Every one of the nine
register values is independent.  The first spare value is arbitrary, and
the remaining spare tail is an arbitrary finite list, possibly empty.
The source paths have length at most 48, so a header distance of at least
64 leaves the required positive delimiter separation.

The intended interface for a nonhalt leaf is a finite positive run to its
successor frame.  Increment changes its chosen logical register by $+1$
and appends a spare value $1$.  Positive decrement changes its register by
$-1$ and increases the extra gap by one.  Zero decrement, break and
continue have their exact source control effects and leave register values
unchanged.  A binary counter rollover shifts the marker and dispatcher
head one position left and adds one to the relative gap.  These motions
are explicitly checked; comparing only a tape viewed relative to the head
would be insufficient.

If $\eta\in\{0,1\}$ is the increment rollover and $\delta=1$ only for a
positive decrement, the nonhalt target has
\[
 h'=h-\eta,\qquad G'=G+\delta+\eta.
 \tag{9}
\]
For continue, $\eta=\delta=0$.  Its target path is the exact stripped path.
All gaps remain nonnegative and no operation empties the spare list.

\subsection{The symbolic rules and why they are universal}

The verification engine represents a tape as homogeneous runs, finite-list
atoms, and an explicit infinite zero exterior.  Run lengths are affine
expressions with nonnegative integer parameters.  It uses only three
execution rules:

\begin{enumerate}
\item \emph{Literal transition.} The currently scanned bit is known and
has positive available length.  The engine applies the exact table row,
including its write, motion, and new state.
\item \emph{Unary self-loop.} If the actual table row loops to the same
state in the traversal direction, induction executes it over a whole
homogeneous run of any natural length.  Zero length is allowed and leaves
the configuration unchanged.
\item \emph{Finite-list traversal.} Before traversing an arbitrary spare
tail, the engine proves the corresponding one-block schema directly on
the exact 238-state table, with arbitrary opaque contexts.  Structural
induction then extends that schema to any finite list.
\end{enumerate}

The four required one-block schemas are shown below.  Each takes exactly
$v+2$ literal steps and moves by the block width in its indicated direction.

\begin{center}
\begin{tabular}{rrll}
\toprule
Entry state & Direction & Input block family & Output family\\
\midrule
73 & $R$ & $B(v)$ & $B(v)$\\
95 & $R$ & $B(v)$ & $C(v)$\\
48 & $L$ & reversed $B(v)$ & reversed $B(v)$\\
112& $L$ & reversed $C(v)$ & reversed $B(v)$\\
\bottomrule
\end{tabular}
\end{center}

Normalization also uses the exact finite-word identity
\[
 C(A)1=1B(A),
 \tag{10}
\]
where $B(A)$ and $C(A)$ concatenate the corresponding blocks of list $A$,
and its reversed version.  For a singleton, (10) is simply
$101^v1=101^{v+1}$.  If it holds for $A$, then
\[
 C(v::A)1=C(v)C(A)1=C(v)1B(A)=1B(v)B(A).
\]
The empty-list case is immediate.  This proves the normalization rule for
every finite spare list, including the empty tail, without a length cutoff.

The generic block proofs reject any access to their opaque exterior.
The halt label has no table rows, so a purported nonhalt proof cannot
silently step through a halt.  List schemas return to their stipulated
nonhalt entry control and have no hidden halting transition.  Unary
shortcuts likewise stay at a working-state self-loop.  These checks
exclude an intermediate halt concealed by a shortcut.

\subsection{Coverage, duration, and iteration}

The final direct replay succeeds at every leaf, splitting each decrement
into the zero and positive cases.  Its exact totals are
\[
 6455+1931=8386\quad\text{guarded cases},
 \qquad 2\,255\,513\quad\text{symbolic proof steps}.
\]
Every one of the 476 literal transition rows is exercised in these proofs.
Actual durations are affine expressions in the gap, nine independent
register parameters, the first spare value, and the remaining list width.
Every coefficient is nonnegative and the smallest constant term is 13.
Thus each certified instruction has a positive finite actual duration
for every concrete valuation.

The engine's limit of 4,000 iterations within one proof attempt bounds
proof construction, not a register value, gap, list length or actual
machine runtime.  Every accepted case has closed before that limit.
The correctness argument is the induction supporting the symbolic rules
and the exact resulting equalities, rather than a long sampled trajectory.

The direct source certificate validates all successor joins, including
(5).  Therefore any successor frame, even one with a shorter program
counter or a zero-only suffix, is exactly the entry tape for its decoded
next leaf.  Equation (9) and spare preservation close the remaining frame
invariants.  Induction over source instructions now gives:

\begin{lemma}[Direct binary implementation]
From any calibrated dispatcher frame of the changed source, the final
$M_{238}$ reaches its halt label if and only if that source execution
reaches its halt instruction.
\end{lemma}

\begin{proof}
A source halt has its own finite certified binary halting segment.
Every nonhalting source instruction has a finite positive binary segment
ending at exactly its required successor frame.  Hence a finite halting
source execution concatenates to a finite binary halt.  An infinite
source execution concatenates infinitely many positive binary segments
with no hidden halt, so the binary execution is infinite.  The source
is deterministic and its full tree always decodes; these two alternatives
also give the reverse implication.
\end{proof}

The primary fresh receipt is \file{evidence/composed238_runtime.json};
the case file records every affine duration.  A separate review in
\file{reviews/runtime_review.md} examines the symbolic rules, the direction
convention, list identities, hidden-halt issue, source normalization and
proof dependencies.  These are executable ordinary mathematical proofs,
not a newly kernel-checked endpoint.

\section{Exact initialization and the ZFC theorem}
\label{sec:boot}

\subsection{A complete finite prefix, checked in two representations}

The initialization target is constructed from the logical seed (6), the
physical permutation, and 167 spare values equal to one.  It is not copied
from an observed trajectory and then used as its own specification.
Two separately coded interpreters execute the final machine: one uses a
sparse set of one-cells and canonical states; the other independently
parses the inherited table into dense transition arrays and uses a dense
byte tape with bounds large enough to contain the entire finite run.

Both produce exactly the following configuration after 409,835 transitions:

\begin{center}
\begin{tabular}{ll}
\toprule
Control & Canonical 84 / inherited 96\\
Head and marker & $h=-5$, marker at $-6$\\
Header & Coordinate 176; distance 181 from the head\\
Extra gap & $181-64=117$\\
Program counter & Empty, equivalently the zero-padded first-leaf path\\
Logical registers & $(3,0,0,0,3,0,0,0,0)$\\
Physical registers & $(0,0,0,3,0,0,3,0,0)$\\
Spare blocks & 167 blocks, each with value 1\\
Tape support & Exactly 351 one-cells; every other cell zero\\
\bottomrule
\end{tabular}
\end{center}

All 351 positions are listed in the new bootstrap receipt.  Both simulations
compare the entire final tape to the independently generated frame, as
well as the state and absolute head coordinate.  The $B_{239}$ predecessor
takes 409,691 steps to the same tape, so this last saving increases this
initialization time by 144 transitions.  The time increase does not add
working states.

\subsection{Composition of the mathematical proof}

Combining exact initialization, the direct binary implementation lemma,
the complete source descriptor equivalence, and the inherited logical
calibration yields
\begin{align*}
 \Halt(M_{238},0^{\mathbb Z})
 &\longleftrightarrow \RHalt(P_1,(\epsilon,R_*))\\
 &\longleftrightarrow \RHalt(P_0,(\epsilon,R_*))\\
 &\longleftrightarrow \operatorname{BProv}(\mathrm{ZFCax\_s},[],\bot)\\
 &\longleftrightarrow \exists p\,\operatorname{Proof}_{\ZFC}(p,\bot).
\end{align*}
This proves (1); negating gives (2).  The new proof uses finite table
facts, induction on natural numbers and finite lists, and finite syntax
translations.  Alongside the particular inherited first-order set-theory
calibration, all these arguments are formalizable in ZFC.  No appeal to
the truth of $\Con(\ZFC)$ occurs in the proof.

The local proof of Section~\ref{sec:local} provides an additional route
through the exact $B_{239}$ predecessor.  Its guard is independently
re-established, so it is not circular with the final-table runtime proof.

\subsection{Consistency, nonhalting, and independence}

The external halting/refutation equivalence established by the ordinary
mathematical proof above implies that, if ZFC is consistent, the machine
actually does not halt.  This semantic conclusion uses that equivalence,
rather than inferring the truth of an arbitrary ZFC theorem from consistency.
Under that same consistency assumption, ZFC cannot prove this nonhalting
statement: such a proof together with (2) would prove $\Con(\ZFC)$,
contrary to G\"odel's second incompleteness theorem.

Bare consistency is not asserted to rule out a ZFC proof of halting.
To conclude two-sided independence, a stronger soundness assumption, such
as arithmetical $\Sigma_1$-soundness, suffices: the false halting assertion
is a $\Sigma_1$ sentence.  This distinction also prevents a misleading
claim of unconditional nonhalting.  The theorem requested here is the
provable equivalence, not a proof of consistency.

The current upstream notes report more proof-assistant progress for the
240-state endpoint than the historical 239 archives did.  Those changing
project statuses do not change the exact tables and source identities
used here.  This submission does not compile a new 238-state endpoint or
claim the separate internal PA result.  Porting the compact local proof
with its selector guard, or instantiating the existing runtime framework
with the new table and source, are concrete next formalization tasks.

\input{formalism.tex}

\section{Reproduction, package contents, and incorporation}
\label{sec:reproduce}

The accepted verification needs Python 3.10 or later with assertions
enabled and only the standard library.  It needs no network, compiler,
SMT solver, Lean installation or original absolute workspace path.
From the extracted directory, run
\begin{verbatim}
python3 -B reproduce.py --output /tmp/zfc238-replay
\end{verbatim}
with a new or empty output directory outside the package.  This replays
the complete final runtime proof as well as construction, source control,
calibrated source-data identification, exact local blocks, bootstrap and
the supplementary finite logical-basis check.
The runtime replay took about 222 seconds in the initial preparation run;
elapsed time is environment dependent.

For a short replay that intentionally omits the unbounded runtime replay,
use
\begin{verbatim}
python3 -B reproduce.py --quick --output /tmp/zfc238-quick
\end{verbatim}
It still binds the archived complete-runtime receipt to the exact table,
source, engine and helper hashes.  Its result explicitly says that the
runtime was not freshly rerun.  To additionally replay the independent
predecessor selector proof, add \file{--with-guard}.  That option performs
a further complete unbounded predecessor run and checks its blank prefix.
No child job is accepted solely for returning process exit code zero;
the driver checks its expected JSON status.

\begin{center}
\begin{tabularx}{\textwidth}{@{}lX@{}}
\toprule
Path & Purpose\\
\midrule
\file{report.tex}, \file{report.pdf} & Complete mathematical report and editable source.\\
\file{machines/} & Exact final tables, bijection and changed source leaves.\\
\file{build.py} & Two exact reconstruction recipes and census.\\
\file{reproduce.py} & Portable replay with honest quick/full status.\\
\file{source/} & Source rotation certificate, leaf arrays, DAG and pinned Lean source data.\\
\file{verification/} & Unbounded runtime engine, exact local certificate, bootstrap and guard checkers.\\
\file{evidence/} & Actual fresh receipts, individual runtime cases and logs.\\
\file{formalism/} & Supplementary logical-basis obstruction and its finite checker.\\
\file{reference/} & Exact predecessor tables, original notices, pins and alternate 43-row recipe.\\
\file{reviews/} & Independent mathematical review of the runtime and proof interfaces.\\
\file{licenses/} & Inherited ProveIt and Wade terms and original notices.\\
\bottomrule
\end{tabularx}
\end{center}

The complete manifest binds all delivered files except itself.  It is an
integrity manifest, not a digital signature.  The packaged quick replay
tests the relocated interfaces; the complete original fresh runtime
receipt remains separately identified.  The inherited proof-assistant
source excerpts are data and reference material, not a standalone Lean
dependency checkout.

A suitable future repository location is
\begin{quote}\small
\file{Computability/BusyBeaver/Research/ZFC/candidate-composed238/}.
\end{quote}
The existing 240-state candidate and both incoming 239-state constructions
should retain their own identities and proof histories.  Incorporation
can import this archive and subsequently add the exact-table kernel
endpoint as a separate, reviewable development.  No remote repository
write was performed during this work.

\section*{Attribution}
The original compact calculus, native set theory, register source and
254-state machine are attributed to Mira-Cult.  The upstream Rocq
verification is attributed to Jeff Ketchersid and contributors; the
compiler framework is attributed to Andrew James Wade.  ProveIt and the
previous incoming submissions retain credit for all prior reductions,
including both 239-state parents used here.  The new work is the concrete
238-state composition, compatibility and guard proofs, direct verification,
new bootstrap checker, supplementary formalism analysis and this report.
No additional restriction is asserted over the newly prepared material;
inherited files retain their supplied terms and notices.

\begin{thebibliography}{20}
\bibitem{rotation240} ProveIt, \emph{Source-rotation 240-state candidate},
inspected at commit\par
\file{ce192939372ac7944fe2a0208c68a96d49fb5e3b}.
\href{https://github.com/VladimirReshetnikov/ProveIt/blob/ce192939372ac7944fe2a0208c68a96d49fb5e3b/Computability/BusyBeaver/Research/ZFC/candidate-rotation240/README.md}{Pinned candidate README}.
\bibitem{incomingA} ProveIt incoming research, \emph{BusyBeaver239 Research},
19 October 2026. \href{https://github.com/VladimirReshetnikov/ProveIt/blob/83fd3585484e5dffb37992f91f7dd4d3e70c437a/docs/incoming/BusyBeaver239_Research_2026-10-19.zip}{Original archive}.
\bibitem{incomingB} ProveIt incoming research, \emph{ZFC239 Companion},
19 October 2026. \href{https://github.com/VladimirReshetnikov/ProveIt/blob/83fd3585484e5dffb37992f91f7dd4d3e70c437a/docs/incoming/ZFC239_Companion_2026-10-19.zip}{Original archive}.
\bibitem{incomingC} ProveIt incoming research, \emph{ZFC 239 State Construction},
19 October 2026. \href{https://github.com/VladimirReshetnikov/ProveIt/blob/83fd3585484e5dffb37992f91f7dd4d3e70c437a/docs/incoming/ZFC_239_State_Construction_2026-10-19.zip}{Original archive}.
\bibitem{calibrated} ProveIt, \file{CalibratedRotation240.lean},
\href{https://github.com/VladimirReshetnikov/ProveIt/blob/13ec39ed43148b1807f51ed6e6528e1c84659562/Computability/BusyBeaver/Lean/BusyBeaver/ZFC/CalibratedRotation240.lean}{pinned source calibration}.
\bibitem{source240} ProveIt, \file{SourceRotation240.lean},
\href{https://github.com/VladimirReshetnikov/ProveIt/blob/13ec39ed43148b1807f51ed6e6528e1c84659562/Computability/BusyBeaver/Lean/BusyBeaver/ZFC/SourceRotation240.lean}{pinned source definition and equivalence}.
\bibitem{seed246} ProveIt, \file{LogicalSeed246.lean},
\href{https://github.com/VladimirReshetnikov/ProveIt/blob/13ec39ed43148b1807f51ed6e6528e1c84659562/Computability/BusyBeaver/Lean/BusyBeaver/ZFC/LogicalSeed246.lean}{pinned seed definition}.
\bibitem{wf} ProveIt, \file{WFTheory.lean},
\href{https://github.com/VladimirReshetnikov/ProveIt/blob/13ec39ed43148b1807f51ed6e6528e1c84659562/Computability/BusyBeaver/Lean/BusyBeaver/ZFC/WFTheory.lean}{native-theory and ZF refutation comparison}.
\bibitem{choice} ProveIt, \file{ConstructibleChoice.lean},
\href{https://github.com/VladimirReshetnikov/ProveIt/blob/13ec39ed43148b1807f51ed6e6528e1c84659562/Computability/BusyBeaver/Lean/BusyBeaver/ZFC/ConstructibleChoice.lean}{ZF and ZFC refutation comparison}.
\bibitem{upstream} Jeff Ketchersid and contributors, \emph{A verified 254-state
Turing machine that halts iff ZFC is inconsistent}, machine by Mira-Cult.
\href{https://github.com/jeff87654/busy-beaver/tree/zf-254/ZF/m254}{Upstream m254 development}.
\bibitem{bbwiki} Busy Beaver Challenge, \emph{Logical independence},
\href{https://wiki.bbchallenge.org/wiki/Logical_independence}{overview},
accessed 19 October 2026.  Used as context; the exact constructions are
identified by their primary artifacts.
\bibitem{friedman} Harvey M. Friedman, \emph{Simplified Axioms for Class
Theory}, 15 September 2020, especially Section 2 and Theorem 3.21.
\href{https://bpb-us-w2.wpmucdn.com/u.osu.edu/dist/1/1952/files/2020/09/NBGAxioms091520.pdf}{Primary manuscript}.
\bibitem{megill} Norman D. Megill, \emph{A Finitely Axiomatized
Formalization of Predicate Calculus with Equality}.
\href{https://us.metamath.org/downloads/finiteaxiom.pdf}{Author's preprint}.
\bibitem{typed} ProveIt, \emph{Typed countermodel},
\href{https://github.com/VladimirReshetnikov/ProveIt/blob/ce192939372ac7944fe2a0208c68a96d49fb5e3b/Computability/BusyBeaver/Research/ZFC/typed_countermodel.md}{pinned equality-free calculus obstruction}.
\end{thebibliography}

\clearpage
\appendix
\section{Complete canonical transition table}
\label{app:table}
The start state is $0$.  The symbol $H$ in this printed table means the
separate halt label $238$.  Each entry $wDq$ means: write bit $w$, move
one cell in direction $D\in\{L,R\}$, and enter $q$.  In the machine file,
all targets are integers and every row is written in the five-field form
\texttt{state read write direction target}.  The printed table is generated
directly from the canonical file bound by the digest in Section~\ref{sec:result}.
\input{transition_table.tex}
\end{document}
